\section{Related work and contribution boundary}

Algorithm selection and instance-space analysis already own the general idea of mapping feature-conditioned performance regions \cite{rice1976algorithm,smithmiles2009meta,smithmiles2023isa}.  Quantum studies already select between quantum and classical algorithms \cite{moussa2020selection}, use instance spaces for quantum optimization \cite{katial2025qaoa}, and predict compiler or device choices \cite{quetschlich2023options,quetschlich2025predictor}.  We do not claim novelty for any of these tasks.

Exact, verified and algebraic circuit optimization is also established.  Prior work supplies verified optimizer transformations \cite{hietala2021verified}, exact pattern matching \cite{iten2022patterns}, graph-theoretic simplification \cite{duncan2020zx}, and optimal Clifford synthesis on bounded qubit counts \cite{bravyi2022clifford}.  Benchmark suites characterize compiler behavior across circuits and devices \cite{li2023qasmbench,quetschlich2023mqtbench}.  These works absorb any claim that exact optimization, formal checking, compiler benchmarking or Clifford synthesis is introduced here.

The Pauli models build on established unitary partitioning and Pauli compilation \cite{izmaylov2020unitary,peres2023pbc,vandenberg2020diagonalization,paykin2023pcoast}.  Connectivity overhead and routing constraints have also been characterized independently \cite{yuan2025connectivity}.  The present structural costs do not replace hardware-aware resource models.

The surviving contribution is a typed combination.  It places constructive gap witnesses, intrinsic versus proof-derived support, objective-conditioned proof validity, representation-level impossibility and prospective falsification in one exact-compiler record.  Each component can succeed or fail independently.  The paper does not introduce a new optimizer, benchmark or algorithm-selection method, and it does not claim that this combination is a universal theory of compiler regimes.
